

\setlength{\parindent}{0pt}
\setlength{\parskip}{6pt plus 2pt minus 1pt}

\newcolumntype{P}[1]{>{\RaggedRight\arraybackslash}p{#1 - 2\tabcolsep - 2\arrayrulewidth}}
\newcolumntype{C}[1]{>{\centering\arraybackslash}p{#1 - 2\tabcolsep - 2\arrayrulewidth}}

\begin{document}

\pagestyle{empty}

\begin{center}
    {\small МИНИСТЕРСТВО НАУКИ И ВЫСШЕГО ОБРАЗОВАНИЯ РОССИЙСКОЙ ФЕДЕРАЦИИ}

    \resizebox{\textwidth}{!}{\fontsize{8pt}{10pt}\selectfont ФЕДЕРАЛЬНОЕ
        ГОСУДАРСТВЕННОЕ АВТОНОМНОЕ ОБРАЗОВАТЕЛЬНОЕ УЧРЕЖДЕНИЕ ВЫСШЕГО ОБРАЗОВАНИЯ}

    {\small \textbf{«МОСКОВСКИЙ ПОЛИТЕХНИЧЕСКИЙ УНИВЕРСИТЕТ»}} \\
    {\small \textbf{(МОСКОВСКИЙ ПОЛИТЕХ)}}
\end{center}

\begin{tabularx}{\textwidth}{@{} p{0.53\textwidth} p{0.43\textwidth} @{}}
     &
    \raggedright
    УТВЕРЖДАЮ                          \\
    заведующая кафедрой                \\
    «Инфокогнитивные технологии»       \\[0.3cm]

    \rule{3cm}{0.4pt} / Е.~А.~Пухова / \\[0.1cm]

    «\rule{0.8cm}{0.4pt}» \rule{3cm}{0.4pt} 2026~г.
\end{tabularx}

\begin{center}
    { \large \textbf{ЗАДАНИЕ}} \\
    {на выполнение курсовой работы (проекта)}

    Дубиной Анастасии Александровне,
\end{center}

обучающейся группы 241-327,

направления подготовки 09.03.01 «Информатика и вычислительная техника»

по дисциплине «Разработка веб-приложений»

на тему «Разработка веб-платформы для публикации, распространения и запуска браузерных игр (WhyGame)»

1. Исходные данные к работе (проекту): документация Flask, PostgreSQL, Unity WebGL, ресурсы сети Интернет, существующие игровые платформы (itch.io, Steam).

2. Содержание задания по курсовой работе (проекту):

{
\renewcommand{\arraystretch}{1.2}
\small
\begin{longtable}{|P{0.5\textwidth}|C{0.11\textwidth}|C{0.15\textwidth}|C{0.24\textwidth}|}

\hline
\textbf{Разрабатываемый вопрос} & \textbf{Объем, \%} & \textbf{Срок} & \textbf{Примечание} \\
\hline
\endhead
\hline
\endfoot

\multicolumn{4}{|P{\textwidth}|}{\textbf{Раздел 1. Анализ предметной области}} \\ \hline
Анализ существующих платформ публикации игр (itch.io, Steam) & 10 & 22.03.2026 & \\ \hline
Анализ технологий веб-разработки и хранения файлов & 7 & 26.03.2026 & \\ \hline
Формулировка цели и задач проекта & 8 & 30.03.2026 & \\ \hline

\multicolumn{4}{|P{\textwidth}|}{\textbf{Раздел 2. Проектирование веб-приложения}} \\ \hline
Анализ целевой аудитории (разработчики и игроки) & 5 & 02.04.2026 & \\ \hline
Описание функциональности (user stories, use-case) & 7 & 06.04.2026 & \\ \hline
Проектирование модели данных (User, Game, Comment, Rating) & 8 & 10.04.2026 & \\ \hline
Проектирование системы хранения игровых файлов & 5 & 14.04.2026 & \\ \hline

\multicolumn{4}{|P{\textwidth}|}{\textbf{Раздел 3. Разработка веб-приложения}} \\ \hline
Разработка архитектуры (Flask + PostgreSQL + файловое хранилище) & 8 & 21.04.2026 & \\ \hline
Реализация аутентификации и авторизации пользователей & 7 & 27.04.2026 & \\ \hline
Реализация системы ролей (разработчик, пользователь, администратор) & 6 & 30.04.2026 & \\ \hline
Реализация CRUD-интерфейса для сущностей (пользователь, игра, комментарий) & 10 & 05.05.2026 & \\ \hline
Реализация загрузки и хранения игровых файлов (Unity WebGL) & 8 & 10.05.2026 & \\ \hline
Реализация запуска игр в браузере & 6 & 14.05.2026 & \\ \hline
Реализация системы комментариев и оценок & 5 & 18.05.2026 & \\ \hline
Агрегирование данных (рейтинги, популярные игры) & 5 & 20.05.2026 & \\ \hline

\multicolumn{4}{|P{\textwidth}|}{\textbf{Раздел 4. Оформление итогов работы}} \\ \hline
Создание Git-репозитория проекта & 3 & 22.03.2026 & \\ \hline
Деплой приложения & 4 & 26.05.2026 & \\ \hline
Оформление отчёта & 3 & 26.05.2026 & \\ \hline

\end{longtable}
}

Руководитель курсовой работы (проекта): \\ старший преподаватель кафедры
«Инфокогнитивные технологии»

{
\renewcommand{\arraystretch}{2}
\noindent
\begin{tabularx}{\textwidth}{@{} P{0.45\textwidth} P{0.3\textwidth} P{0.25\textwidth} @{}}

    «\rule{0.8cm}{0.4pt}» \quad \rule{2cm}{0.4pt} \quad 2026~г. &
    \centering \rule{3.5cm}{0.4pt} \par \footnotesize (подпись) &
    \hfill А.~С.~Кружалов                                         \\
\end{tabularx}

\begin{tabularx}{\textwidth}{@{} P{0.45\textwidth} P{0.55\textwidth} @{}}
    Дата выдачи задания           &
    \hfill «30» \quad марта \quad 2026~г. \\

    Дата сдачи выполненной работы &
    \hfill «\rule{0.8cm}{0.4pt}» \quad \rule{2cm}{0.4pt} \quad 2026~г. \\
\end{tabularx}
}

Задание приняла к исполнению

{
\renewcommand{\arraystretch}{1.5}
\begin{tabularx}{\textwidth}{@{} P{0.45\textwidth} P{0.3\textwidth} P{0.25\textwidth} @{}}

     \hfill «30» \quad марта \quad 2026~г. &
    \centering \rule{3.5cm}{0.4pt} \par \footnotesize (подпись) &
    \hfill А. А. Дубина                                           \\
\end{tabularx}
}

\end{document}